\chapter{Espectro, resolvente y modelos concretos}\label{cap:v3-05}
En dimensión finita, resolver $(T-\lambda I)x=y$ equivale a evitar los valores propios. En dimensión infinita aparecen dos dificultades adicionales: la imagen puede ser densa sin ser todo el espacio, o puede dejar fuera un subespacio detectable por funcionales. La teoría espectral organiza estas posibilidades. En todo el capítulo $X$ será un espacio de Banach complejo \emph{no nulo}; escribiremos $\lambda$ por $\lambda I$ cuando no haya ambigüedad.

\section{Invertibilidad y tipos de espectro}\label{lec:v3-07}
\fuentelectura{7}{El espectro de un operador}
\begin{lema}[Serie de Neumann]\label{v3:neumann}\index{Neumann, serie de}
Si $A\in\Bcal(X)$ y $\norm A<1$, entonces
\[(I-A)^{-1}=\sum_{n=0}^{\infty}A^n,\qquad
\norm{(I-A)^{-1}}\le\frac1{1-\norm A}.
\]
El conjunto de operadores invertibles es abierto. Si $T$ es invertible y $\norm{S-T}<\norm{T^{-1}}^{-1}$, entonces $S$ es invertible.
\end{lema}
\begin{proof}
La serie converge en el Banach $\Bcal(X)$ porque $\sum\norm{A^n}\le\sum\norm A^n<\infty$. Para $P_n=\sum_{k=0}^nA^k$ se tiene $(I-A)P_n=P_n(I-A)=I-A^{n+1}$. La continuidad del producto y $A^{n+1}\to0$ prueban ambas identidades de inversa; la suma geométrica da la cota. Para la última afirmación, escribimos
\[S=T[I+T^{-1}(S-T)]
\]
y aplicamos lo anterior al segundo factor. El producto de dos operadores invertibles tiene inversa en orden inverso: $(ST)(T^{-1}S^{-1})=I=(T^{-1}S^{-1})(ST)$. Esta verificación también resuelve la demostración propuesta en la proposición 2 original. Finalmente, $1=\norm I\le\norm T\norm{T^{-1}}$.
\end{proof}

\begin{definicion}\index{espectro}\index{resolvente}\index{espectro!puntual}\index{espectro!continuo}\index{espectro!residual}
El \emph{resolvente} de $T$ es $\rho(T)=\{\lambda:T-\lambda I\text{ es invertible}\}$ y su \emph{espectro} es $\sigma(T)=\C\setminus\rho(T)$. Usaremos $r(T)$ para el radio espectral, evitando confundirlo con el conjunto $\rho(T)$. El espectro se divide en tres partes disjuntas:
\begin{itemize}
\item $\sigma_p(T)$: $T-\lambda I$ no es inyectivo; sus elementos son valores propios.
\item $\sigma_c(T)$: $T-\lambda I$ es inyectivo, tiene imagen densa y no es suprayectivo.
\item $\sigma_r(T)$: $T-\lambda I$ es inyectivo y su imagen no es densa.
\end{itemize}
\end{definicion}
Si $T-\lambda I$ fuese biyectivo, su inversa sería acotada por aplicación abierta. Por eso las tres alternativas anteriores agotan el espectro. En dimensión finita inyectividad equivale a suprayectividad; allí todo el espectro es puntual.

\begin{ejemplo}[Multiplicación por la coordenada]\label{v3:mult-coordenada}
En $L^1[0,1]$, $Mf(t)=tf(t)$ tiene norma uno y
\[\sigma(M)=\sigma_c(M)=[0,1].
\]
La cota $\norm{Mf}_1\le\norm f_1$ es inmediata; los indicadores normalizados de intervalos próximos a uno muestran que la norma es uno. Si $\lambda\notin[0,1]$, la función $(t-\lambda)^{-1}$ es acotada y multiplicar por ella da la inversa. Para cualquier $\lambda$, la igualdad $(t-\lambda)f(t)=0$ casi en todas partes implica $f=0$, pues el posible conjunto excepcional $\{\lambda\}\cap[0,1]$ tiene medida cero. Si $\lambda\in[0,1]$, para $g\in L^1$ definimos
\[f_\delta(t)=\frac{g(t)}{t-\lambda}\boldsymbol1_{\{|t-\lambda|\ge\delta\}}.
\]
Es integrable porque $|f_\delta|\le|g|/\delta$, y $(M-\lambda)f_\delta\to g$ en $L^1$ por convergencia dominada. La imagen es densa. No contiene la función constante uno: su única preimagen formal sería $1/(t-\lambda)$, no integrable incluso si $\lambda$ es un extremo. Así la imagen es propia. Esta prueba desarrolla la aproximación propuesta en la lectura y evita los errores de signos de sus intervalos de corte.
\end{ejemplo}

\begin{ejemplo}[Corrimiento derecho en $\ell^\infty$]\label{v3:shift-infinito}\index{corrimiento!derecho}
Sea $R(x_1,x_2,\ldots)=(0,x_1,x_2,\ldots)$. Entonces $\norm R=1$ y
\[\sigma(R)=\sigma_r(R)=\{\lambda:|\lambda|\le1\}.
\]
El operador $R-\lambda$ es inyectivo: para $\lambda=0$ esto es directo; para $\lambda\ne0$, la primera coordenada y luego las siguientes fuerzan sucesivamente $x_n=0$. Fuera del disco unidad es invertible por Neumann. Si $|\lambda|<1$, la funcional no nula
\[\phi(y)=\sum_{n\ge1}\lambda^{n-1}y_n
\]
es continua y se anula en $\ran(R-\lambda)$, como se comprueba desplazando el índice de una serie absolutamente convergente. Por tanto esa imagen no es densa.

Si $|\lambda|=1$, la ecuación $(R-\lambda)x=y$ determina por inducción
\[x_n=-\lambda^{-n}\sum_{k=1}^n\lambda^{k-1}y_k.
\]
Tomemos $y_k=\lambda^{1-k}$. Si $\norm z_\infty<1/2$ y $(R-\lambda)x=y+z$, entonces
\[|x_n|=\left|n+\sum_{k=1}^n\lambda^{k-1}z_k\right|
\ge n-n\norm z_\infty>n/2.
\]
No existe tal $x\in\ell^\infty$, y la bola $B(y,1/2)$ no intersecta la imagen. Esto prueba también en la frontera el carácter residual. Los exponentes y el signo de esta recurrencia corrigen las fórmulas de la lectura original.
\end{ejemplo}

\section{Compacidad del espectro y cálculo polinómico}\label{lec:v3-08}
\fuentelectura{8}{Propiedades del espectro de un operador}
\begin{teorema}\label{v3:espectro-compacto}
El espectro de todo $T\in\Bcal(X)$ es un compacto no vacío, contenido en $\{z:|z|\le\norm T\}$.
\end{teorema}
\begin{proof}
La apertura de los invertibles implica que $\rho(T)$ es abierto. Si $|z|>\norm T$, la factorización $zI-T=z(I-T/z)$ y Neumann dan invertibilidad y
\[\norm{(zI-T)^{-1}}\le\frac1{|z|-\norm T}.
\]
Así el espectro es cerrado y acotado, por tanto compacto. Para la no vacuidad, supongamos que fuera vacío y pongamos $R(z)=(zI-T)^{-1}$ para todo $z$. La identidad de inversas da
\[R(z)-R(w)=(w-z)R(z)R(w).
\]
Además, la expansión de Neumann alrededor de cualquier $w$ prueba que $R(z)\to R(w)$ en norma. Dividiendo la identidad anterior por $z-w$ obtenemos $R'(w)=-R(w)^2$. Para cada $\Phi\in\Bcal(X)^*$, la función escalar $f(z)=\Phi(R(z))$ es entera. La cota anterior la acota fuera de un disco, y la continuidad la acota dentro. Por Liouville es constante, y como $f(z)\to0$ cuando $|z|\to\infty$, es idénticamente cero. Hahn--Banach, aplicado al espacio $\Bcal(X)$, implica $R(z)=0$, imposible porque $R(z)(zI-T)=I\ne0$. La hipótesis $X\ne\{0\}$ es la que asegura esta última desigualdad.
\end{proof}
La demostración sólo usa el teorema escalar de Liouville, que se deduce de la fórmula de Cauchy: para una función entera acotada por $M$, $|f'(z)|\le M/R$ en circunferencias de radio arbitrario $R$, luego $f'=0$.

\begin{lema}\label{v3:producto-conmutativo}
Un producto finito de operadores que conmutan dos a dos es invertible si y sólo si cada factor lo es.
\end{lema}
\begin{proof}
Basta tratar dos factores $A,B$. Una dirección usa $(AB)^{-1}=B^{-1}A^{-1}$. Si $AB$ es invertible, $A$ y $B$ conmutan con $AB$, luego también con $(AB)^{-1}$: de $AC=CA$ se deduce $C^{-1}A=AC^{-1}$ multiplicando por $C^{-1}$ a ambos lados. Así $B(AB)^{-1}$ es a la vez inversa derecha e izquierda de $A$; análogamente $A(AB)^{-1}$ lo es de $B$. La inducción da el caso finito.
\end{proof}

\begin{teorema}[Aplicación espectral polinómica]\label{v3:espectral-polinomios}\index{aplicación espectral}
Para todo polinomio complejo $p$, $\sigma(p(T))=p(\sigma(T))$.
\end{teorema}
\begin{proof}
Si $p=c$ es constante, $p(T)=cI$ tiene espectro $\{c\}$, igual a $p(\sigma(T))$ por no vacuidad. Si $p$ tiene grado $m\ge1$, para cada $\mu$ el teorema fundamental del álgebra da
\[p(z)-\mu=a\prod_{j=1}^m(z-\lambda_j),\qquad a\ne0.
\]
Sustituyendo $T$, los factores conmutan. El lema anterior dice que $p(T)-\mu$ es invertible exactamente cuando todos los $T-\lambda_j$ lo son, es decir, cuando ningún punto de $\sigma(T)$ satisface $p(\lambda)=\mu$. Esto equivale a $\mu\notin p(\sigma(T))$ y prueba la igualdad.
\end{proof}

\begin{teorema}[Espectro y dualidad]\label{v3:espectro-adjunto}
Para el adjunto de Banach se cumple
\[\sigma(T')=\sigma(T),\qquad
\sigma_r(T)\subset\sigma_p(T')\subset\sigma_p(T)\cup\sigma_r(T),\qquad
\sigma_c(T')\subset\sigma_c(T).
\]
\end{teorema}
\begin{proof}
Si $A$ es invertible, $A'$ tiene inversa $(A^{-1})'$ por la regla de composición. Si $A'$ es invertible, su imagen densa implica que $A$ es inyectiva, y su núcleo cero implica que $\ran A$ es densa, por el ejercicio~\ref{ej:v3-l04-5}. Además $\ran A'$ es cerrada, así que $\ran A$ es cerrada por el teorema~\ref{v3:rango-cerrado}. Es, pues, todo $X$. Aplicando esto a $A=T-\lambda$ y usando $A'=T'-\lambda$ se obtiene la igualdad de espectros.

Si $\lambda\in\sigma_r(T)$, la imagen no densa de $T-\lambda$ tiene un anulador no nulo; éste es $\ker(T'-\lambda)$, luego $\lambda\in\sigma_p(T')$. A la inversa, un vector no nulo de ese núcleo demuestra que la imagen de $T-\lambda$ no es densa. Por tanto $\lambda$ no puede estar en $\sigma_c(T)$ ni en el resolvente, y pertenece a $\sigma_p(T)\cup\sigma_r(T)$. Finalmente, si $\lambda\in\sigma_c(T')$, la inyectividad de $T'-\lambda$ implica densidad de la imagen de $T-\lambda$, y la densidad de la imagen del primero implica inyectividad del segundo. La igualdad de espectros excluye la suprayectividad, de modo que $\lambda\in\sigma_c(T)$. No se ha confundido el anulador cero con densidad en norma de un subespacio del dual: ese paso, empleado sin hipótesis en la prueba original, se sustituye aquí por rango cerrado.
\end{proof}

\section{Ejercicios de la lectura 7}
\begin{ejercicio}\label{ej:v3-l07-1}
Sea $(a_n)$ acotada y $Dx=(a_nx_n)$ en $\ell^1$. Calcule su espectro.
\end{ejercicio}


\begin{ejercicio}\label{ej:v3-l07-2}
Demuestre que el corrimiento derecho bilateral en $\ell^2(\Z)$ es invertible y determine su inversa.
\end{ejercicio}


\begin{ejercicio}\label{ej:v3-l07-3}
Para $f\in C[0,1]$, sea $Tg(t)=f(1-t)g(t)$ en $L^1[0,1]$. Pruebe que $\sigma(T)=f([0,1])$ y determine los valores propios.
\end{ejercicio}


\begin{ejercicio}\label{ej:v3-l07-4}
Los vectores propios asociados a valores propios \emph{distintos} $\lambda_1,\ldots,\lambda_m$ son linealmente independientes.
\end{ejercicio}


\begin{ejercicio}\label{ej:v3-l07-5}
Si $T\in\Bcal(X)$ es invertible, pruebe $\sigma(T^{-1})=\{\lambda^{-1}:\lambda\in\sigma(T)\}$.
\end{ejercicio}


\begin{ejercicio}\label{ej:v3-l07-6}
Para el adjunto de Hilbert, pruebe $\sigma(T^*)=\{\cl\lambda:\lambda\in\sigma(T)\}$.
\end{ejercicio}


\begin{ejercicio}\label{ej:v3-l07-7}
Para $L(x_1,x_2,\ldots)=(x_2,x_3,\ldots)$ en $\ell^1$, pruebe (i) $\sigma_p(L)=\{|z|<1\}$; (ii) $\sigma_c(L)=\{|z|=1\}$; (iii) $\sigma_r(L)=\varnothing$.
\end{ejercicio}


\section{Ejercicios de la lectura 8}
\begin{ejercicio}\label{ej:v3-l08-1}
Demuestre que un producto finito de operadores que conmutan dos a dos es invertible si y sólo si cada factor lo es.
\end{ejercicio}


\begin{ejercicio}\label{ej:v3-l08-2}
Pruebe: (i) $T^2=T$ implica $\sigma(T)\subset\{0,1\}$; (ii) $T$ nilpotente implica $\sigma(T)=\{0\}$.
\end{ejercicio}


\begin{ejercicio}\label{ej:v3-l08-3}
Si $T_n\to T$ en norma, $\lambda_n\to\lambda$ y $\lambda_n\in\sigma(T_n)$, pruebe $\lambda\in\sigma(T)$.
\end{ejercicio}


\begin{ejercicio}\label{ej:v3-l08-4}
Para los corrimientos unilaterales $L,R$ en $\ell^\infty$, demuestre $\sigma(LR)\ne\sigma(RL)$.
\end{ejercicio}


\begin{ejercicio}\label{ej:v3-l08-5}
Si $(S\times T)(x,y)=(Sx,Ty)$, pruebe $\sigma(S\times T)=\sigma(S)\cup\sigma(T)$.
\end{ejercicio}


\begin{ejercicio}\label{ej:v3-l08-6}
Si $T_n\to T$ en norma, para todo $\eps>0$ existe $N$ con $\sigma(T_n)\subset\sigma(T)+B(0,\eps)$ cuando $n\ge N$.
\end{ejercicio}


\begin{ejercicio}\label{ej:v3-l08-7}
Todo compacto no vacío $K\subset\C$ es el espectro de un operador acotado en un espacio de Banach.
\end{ejercicio}


\begin{ejercicio}\label{ej:v3-l08-8}
Los operadores acotados inyectivos con imagen cerrada forman un abierto de $\Bcal(X)$.
\end{ejercicio}


\section{Radio espectral y crecimiento de potencias}\label{lec:v3-09}
\fuentelectura{9}{El radio espectral}
\begin{definicion}\index{radio espectral}
Definimos $r(T)=\max\{|\lambda|:\lambda\in\sigma(T)\}$. El máximo existe por compacidad y $0\le r(T)\le\norm T$. Un operador con $r(T)=0$ se llama \emph{cuasinilpotente}.
\end{definicion}
Un nilpotente no nulo muestra que la desigualdad con la norma puede ser estricta: en los polinomios de grado a lo sumo $m$, la derivada tiene potencia $m+1$ igual a cero, aunque no es el operador cero si $m\ge1$.

\begin{teorema}[Fórmula de Gelfand]\label{v3:gelfand-radio}\index{Gelfand!fórmula del radio}
Para $T\in\Bcal(X)$,
\[r(T)=\lim_{n\to\infty}\norm{T^n}^{1/n}.
\]
\end{teorema}
\begin{proof}
Aplicación espectral da $\lambda^n\in\sigma(T^n)$ para $\lambda\in\sigma(T)$; luego $r(T)\le\norm{T^n}^{1/n}$ para todo $n$. Falta la cota para el límite superior.

Sea $D=\{z:|z|<1/r(T)\}$, interpretado como $\C$ si $r(T)=0$. Para $z\ne0$ en $D$, $1/z\notin\sigma(T)$ y $I-zT$ es invertible. En $z=0$ también lo es. El mismo argumento local de Neumann usado para el resolvente muestra que $z\mapsto(I-zT)^{-1}$ es holomorfa en norma en $D$. Para cada $\Phi\in\Bcal(X)^*$, la función escalar
\[g_\Phi(z)=\Phi((I-zT)^{-1})
\]
es holomorfa allí. Cerca de cero, Neumann da $g_\Phi(z)=\sum_{n\ge0}\Phi(T^n)z^n$. Por unicidad de los coeficientes de Taylor y el teorema de expansión de una función holomorfa en un disco, esta serie converge para todo $z\in D$. Este último hecho se obtiene de la fórmula de Cauchy en cualquier circunferencia compactamente contenida en $D$; no se presupone convergencia de una serie de operadores fuera del disco $|z|<1/\norm T$.

Fijemos $z\in D\setminus\{0\}$. La convergencia escalar implica que $(z^n\Phi(T^n))_n$ es acotada para cada $\Phi$. En el Banach $\Bcal(X)^*$ consideramos las funcionales $F_n(\Phi)=\Phi(z^nT^n)$, cuya norma es $\norm{z^nT^n}$ por la inserción isométrica en el bidual. Acotación uniforme proporciona $C_z<\infty$ con $|z|^n\norm{T^n}\le C_z$. Por tanto
\[\limsup_n\norm{T^n}^{1/n}\le |z|^{-1}.
\]
Si $r(T)>0$, hacemos $|z|\uparrow1/r(T)$; si $r(T)=0$, hacemos $|z|\to\infty$. En ambos casos el límite superior es a lo sumo $r(T)$, y la cota inicial concluye. Aquí las funcionales actúan sobre $\Bcal(X)$, no sobre $X$; además el argumento correcto es $I-zT$, dos precisiones necesarias en la prueba original.
\end{proof}

\begin{ejemplo}[Cuasinilpotente que no es nilpotente]
En $\ell^\infty$, sea $Tx=(0,x_1,x_2/2,x_3/3,\ldots)$. Por inducción,
\[(T^nx)_k=\begin{cases}0,&k\le n,\\
\displaystyle\frac{x_{k-n}}{(k-1)(k-2)\cdots(k-n)},&k>n.
\end{cases}
\]
El denominador mínimo es $n!$, alcanzado cuando $k=n+1$. Así $\norm{T^n}\le1/n!$ y la igualdad se alcanza en $e_1$, pues $T^ne_1=e_{n+1}/n!$. Además $(n!)^{1/n}\to\infty$: al menos $\lfloor n/2\rfloor$ factores son mayores o iguales que $n/2$. La fórmula de Gelfand da $r(T)=0$, de modo que $\sigma(T)=\{0\}$. El operador es inyectivo y su imagen tiene primera coordenada cero, por tanto cero es residual. Ninguna potencia es cero, como muestra la misma fórmula en $e_1$.
\end{ejemplo}

\section{Ejercicios de la lectura 9}
\begin{ejercicio}\label{ej:v3-l09-1}
Si $ST=TS$, pruebe $r(ST)\le r(S)r(T)$ y $r(S+T)\le r(S)+r(T)$.
\end{ejercicio}


\begin{ejercicio}\label{ej:v3-l09-2}
Pruebe $r(T^m)=r(T)^m$ para todo entero $m\ge1$.
\end{ejercicio}


\begin{ejercicio}\label{ej:v3-l09-3}
Para una matriz compleja $A$ con la norma de operador euclídea, pruebe $\norm A^2=r(A^*A)$.
\end{ejercicio}


\begin{ejercicio}\label{ej:v3-l09-4}
Demuestre que $T\mapsto r(T)$ es semicontinua superior para la norma de operadores.
\end{ejercicio}


\section{Ampliación: multiplicadores y corrimientos bilaterales}
La siguiente formulación desarrolla el modelo presentado por \textbf{Gerardo Ramos Vázquez}, \emph{Los operadores de multiplicación en espacios de medida $\sigma$-finita} \cite{ramos2018}, alojado en el sitio de Egor Maximenko. La atribución corresponde al autor de esa presentación; las demostraciones y ejercicios siguientes son desarrollos de esta edición.

\begin{teorema}\label{v3:multiplicadores}\index{operador!de multiplicación}\index{rango esencial}
Sea $(\Omega,\mu)$ un espacio de medida $\sigma$-finito con $L^2(\mu)\ne\{0\}$ y $a\in L^\infty(\mu)$. El operador $M_af=af$ satisface
\[\norm{M_a}=\norm a_\infty,\quad M_a^*=M_{\cl a},\quad
\sigma(M_a)=\operatorname{essran}(a),
\]
donde $\lambda\in\operatorname{essran}(a)$ significa que $\mu\{|a-\lambda|<\eps\}>0$ para todo $\eps>0$. Sus valores propios son los $\lambda$ con $\mu\{a=\lambda\}>0$; el resto del espectro es continuo y el residual es vacío.
\end{teorema}
\begin{proof}
La desigualdad $\norm{af}_2\le\norm a_\infty\norm f_2$ da una cota. Si $0<c<\norm a_\infty$, el conjunto $\{|a|>c\}$ tiene medida positiva. Por $\sigma$-finitud contiene un subconjunto $E$ de medida positiva finita. El vector $f=\boldsymbol1_E/\sqrt{\mu(E)}$ tiene norma uno y $\norm{M_af}_2\ge c$. Al hacer $c\uparrow\norm a_\infty$ se obtiene la igualdad; si la norma es cero no hace falta este paso. La fórmula del adjunto se deduce de
$\int af\cl g\,d\mu=\int f\cl{\cl a g}\,d\mu$.

Si $\lambda$ no está en el rango esencial, $|a-\lambda|\ge\eps$ casi en todas partes para algún $\eps>0$, y $M_{1/(a-\lambda)}$ es una inversa acotada. Si está, los conjuntos $\{|a-\lambda|<1/n\}$ contienen subconjuntos $E_n$ de medida positiva finita. Sus indicadores normalizados $f_n$ satisfacen $\norm{(M_a-\lambda)f_n}_2\le1/n$, lo que excluye una inversa acotada. Esto prueba la igualdad de espectros.

La ecuación $(a-\lambda)f=0$ equivale a que $f$ esté soportada, salvo conjuntos nulos, en $\{a=\lambda\}$. Por $\sigma$-finitud existe allí un indicador integrable no nulo exactamente cuando ese conjunto tiene medida positiva. Si su medida es cero, el operador es inyectivo; además $\ker(M_a-\lambda)^*=\ker M_{\cl a-\cl\lambda}=0$, de modo que su imagen es densa. Si $\lambda$ está en el espectro, esa imagen no es todo el espacio y el punto es continuo. Por ello nunca aparece espectro residual.
\end{proof}

\begin{ejercicio}[Complemento: norma y espectro de una proyección]\label{ej:v3-ampl-multiplicador}
Sea $E\subset\Omega$ medible y suponga que $E$ y su complemento tienen medida positiva. Determine norma, espectro y espacios propios de $M_{\boldsymbol1_E}$ en $L^2(\mu)$.
\end{ejercicio}


\begin{ejercicio}[Complemento: espectro del corrimiento bilateral]\label{ej:v3-ampl-shift}
Para $R$ en $\ell^2(\Z)$ del ejercicio~\ref{ej:v3-l07-2}, pruebe $\sigma(R)=\sigma_c(R)=\{|z|=1\}$.
\end{ejercicio}


Los modelos de multiplicación y corrimiento muestran que conocer el conjunto espectral no basta para conocer el operador: el mismo disco puede tener espectro puntual o residual según el espacio y la dirección del corrimiento. La fórmula de Gelfand añade una información cuantitativa, el crecimiento asintótico de las potencias. En el siguiente capítulo la compacidad reducirá gran parte de la teoría espectral a espacios propios de dimensión finita.
